\documentclass[10pt]{article}
\usepackage[a4paper,margin=0.6in]{geometry}
\usepackage{amsmath,amssymb,amsfonts,mathtools}
\allowdisplaybreaks
\setlength{\jot}{5pt}
\numberwithin{equation}{section}

\title{Group I, Row III: explicit step-by-step simplification}
\author{}
\date{}

\begin{document}
\maketitle
\tableofcontents

\section{Result in one place}

The four pieces $N_{1},N_{2},N_{3a},N_{3b}$ of Group I, Row III contain 22 terms. They reduce to
\begin{align}
\frac{d^{3}N_{\rm I.III}}{d^{3}k}&=\frac{1}{(2\pi)^{3}}\,\frac{i g^{4}}{8\pi^{5}}\,\frac{1}{k^{+}}\int_{w,w',x,y,z,x'}e^{-ik\cdot(w'-w)}\;
\frac{(x'-w')^{i}\,(y-z)^{j}}{(x'-w')^{2}\,(y-z)^{2}}\;\mathcal{R}^{ij}\;\mathcal{W}\;+\;{\rm c.c.},\\[4pt]
\mathcal{R}^{ij}&=K^{ij}\left\{\log\frac{(x-z)^{2}}{(x-w)^{2}}+i\,\sigma\,\pi\right\}
-\frac12\,S^{ij}\,\log\frac{(z-w)^{2}}{(x-w)^{2}}+\mathcal{A}_{0}\,T^{ij}, \label{eq:main}
\end{align}
with the colour factor $\mathcal{W}$ of \eqref{eq:W}, the tensors $K,S,T$ of \eqref{eq:KST} and $\sigma=\pm1$ fixed by the $i\epsilon$ prescription. In two transverse dimensions this becomes \eqref{eq:2D}, and the last two terms are a single one-parameter integral, \eqref{eq:irep}.

\section{Input and notation}

\subsection{What is common to all 22 terms}
Every term of $N_{1},N_{2},N_{3a},N_{3b}$ contains
\begin{itemize}
\item the measure and phase $\dfrac{1}{(2\pi)^{3}}\displaystyle\int_{w,w'}e^{-ik\cdot(w'-w)}\int_{x,y,z,x'}$ and the factor $\dfrac{1}{k^{+}}$;
\item the prefactor $\dfrac{ig^{4}f^{abd}}{8\pi^{5}}$, or $\dfrac{ig^{4}f^{abd}}{16\pi^{5}}=\dfrac12\cdot\dfrac{ig^{4}f^{abd}}{8\pi^{5}}$;
\item the two vectors $\dfrac{(x'-w')}{(x'-w')^{2}}$ and $\dfrac{(y-z)}{(y-z)^{2}}$, each contracted with something;
\item the same Wilson-line factor
\begin{align}
\mathcal{W}\equiv\Big[U^{ab'}(x')-U^{ab'}(w')\Big]\rho^{b'}(x')\;f^{abd}\,U^{de}(x)\Big[U^{bc}(z)-U^{bc}(y)\Big]\rho^{e}(x)\,\rho^{c}(y). \label{eq:W}
\end{align}
\end{itemize}
$\mathcal{W}$ does not depend on $w$.

\subsection{Notation}
\begin{align}
&a\equiv x-w,\qquad b\equiv x-z,\qquad c\equiv z-w=a-b,\qquad A\equiv a^{2},\quad B\equiv b^{2},\quad C\equiv c^{2}=A+B-2\,a\cdot b,\notag\\
&u^{i}\equiv\frac{(x'-w')^{i}}{(x'-w')^{2}},\qquad v^{j}\equiv\frac{(y-z)^{j}}{(y-z)^{2}} .
\end{align}
Identities used below:
\begin{align}
c\cdot a=A-a\cdot b,\qquad c\cdot b=a\cdot b-B,\qquad A+B-C=2\,a\cdot b. \label{eq:ids}
\end{align}
The four logarithms and the function $\mathcal{A}_{0}$ that occur are
\begin{align}
L_{a}&\equiv\log\frac{B}{B-A}=\log\frac{-B}{A-B}, & L_{b}&\equiv\log\frac{A-B}{A}, & L_{c}&\equiv\log\frac{-B}{A}, & L_{d}&\equiv\log\frac{C}{A},
\end{align}
\begin{align}
\mathcal{A}_{0}=\frac{1}{\Delta}\left\{\arctan\frac{B-a\cdot b}{\Delta}+\arctan\frac{a\cdot b}{\Delta}\right\},\qquad\Delta\equiv\sqrt{AB-(a\cdot b)^{2}} .
\end{align}
$L_{a}$ is the logarithm of the first term of $N_{1}$ and of all six terms of $N_{2}$ (written there as $\log[-(x-z)^{2}/((x-w)^{2}-(x-z)^{2})]$). $L_{b}$ is the logarithm of the second term of $N_{1}$ and of the first six terms of $N_{3a}$. $L_{c}$ occurs once, in the third term of $N_{1}$. $L_{d}$ occurs in the last three terms of $N_{3b}$.

\subsection{Definition of the kernel}
Writing each term as $u^{i}v^{j}\times(\text{tensor})^{ij}$, Row III is
\begin{align}
\frac{d^{3}N_{\rm I.III}}{d^{3}k}=\frac{1}{(2\pi)^{3}}\,\frac{i g^{4}}{8\pi^{5}}\,\frac{1}{k^{+}}\int_{w,w',x,y,z,x'}e^{-ik\cdot(w'-w)}\;u^{i}\,v^{j}\;\mathcal{R}^{ij}\;\mathcal{W}+{\rm c.c.},\qquad
\mathcal{R}^{ij}=\sum_{n=1}^{22}t_{n}^{ij}. \label{eq:def}
\end{align}
Example of the translation: the first term of $N_{2}$ is
\begin{align}
-\frac{ig^{4}f^{abd}}{8\pi^{5}}\,\log\!\left[\frac{-(x-z)^{2}}{(x-w)^{2}-(x-z)^{2}}\right]\frac{(x'-w')\cdot(z-w)\;(y-z)\cdot(x-w)}{(x'-w')^{2}\,(y-z)^{2}\,(z-w)^{2}\left[(x-w)^{2}-(x-z)^{2}\right]}
\;\;\longrightarrow\;\;t^{ij}=-L_{a}\,\frac{c^{i}a^{j}}{C\,(A-B)},
\end{align}
where the overall minus sign is the sign in front of the braces of $N_{2}$.

\section{Step 1: the 22 terms}

In units of the prefactor of \eqref{eq:def}:

\paragraph{$N_{1}$ (three terms).}
\begin{align}
t_{1}&=+\frac12\,\delta^{ij}\,\frac{A+B}{(B-A)^{2}}\,L_{a}, &
t_{2}&=+\frac12\,\delta^{ij}\,\frac{A+B}{(B-A)^{2}}\,L_{b}, &
t_{3}&=-\frac{a^{i}b^{j}}{AB}\,L_{c}.
\end{align}

\paragraph{$N_{2}$ (six terms, overall sign $-$, all with $L_{a}$).}
\begin{align}
t_{4}&=-L_{a}\,\frac{c^{i}a^{j}}{C\,(A-B)}, &
t_{5}&=+L_{a}\,\frac{B}{A}\,\frac{a^{i}c^{j}}{C\,(A-B)}, &
t_{6}&=-L_{a}\,\delta^{ij}\,\frac{(c\cdot a)\,B}{C\,(A-B)^{2}},\notag\\
t_{7}&=+L_{a}\,\frac{A}{B}\,\frac{c^{i}b^{j}}{C\,(A-B)}, &
t_{8}&=-L_{a}\,\frac{b^{i}c^{j}}{C\,(A-B)}, &
t_{9}&=+L_{a}\,\delta^{ij}\,\frac{(c\cdot b)\,A}{C\,(A-B)^{2}}.
\end{align}

\paragraph{$N_{3a}$ (nine terms).}
\begin{align}
t_{10}&=-L_{b}\,\delta^{ij}\,\frac{(c\cdot a)\,B}{C\,(A-B)^{2}}, &
t_{11}&=+L_{b}\,\delta^{ij}\,\frac{(c\cdot b)\,A}{C\,(A-B)^{2}}, &
t_{12}&=-L_{b}\,\frac{c^{i}a^{j}}{C\,(A-B)},\notag\\
t_{13}&=+L_{b}\,\frac{A}{B}\,\frac{c^{i}b^{j}}{C\,(A-B)}, &
t_{14}&=+L_{b}\,\frac{B}{A}\,\frac{a^{i}c^{j}}{C\,(A-B)}, &
t_{15}&=-L_{b}\,\frac{b^{i}c^{j}}{C\,(A-B)},\notag\\
t_{16}&=+\mathcal{A}_{0}\,\frac{c^{i}a^{j}}{C}, &
t_{17}&=+\mathcal{A}_{0}\,\frac{b^{i}c^{j}}{C}, &
t_{18}&=-\mathcal{A}_{0}\,\frac{(a\cdot b)\,c^{i}b^{j}}{C\,B}.
\end{align}

\paragraph{$N_{3b}$ (four terms).}
\begin{align}
t_{19}&=-\mathcal{A}_{0}\,\frac{(a\cdot b)\,a^{i}c^{j}}{C\,A}, &
t_{20}&=-\frac12\,L_{d}\,\frac{c^{i}b^{j}}{C\,B}, &
t_{21}&=+\frac12\,L_{d}\,\frac{a^{i}c^{j}}{C\,A}, &
t_{22}&=-\frac12\,L_{d}\,\frac{\delta^{ij}}{C}.
\end{align}

\section{Step 2: group by logarithm}

\begin{center}
\renewcommand{\arraystretch}{1.3}
\begin{tabular}{c|l}
function & terms\\ \hline
$L_{a}$ & $t_{1},\ t_{4},\dots,t_{9}$\\
$L_{b}$ & $t_{2},\ t_{10},\dots,t_{15}$\\
$L_{c}$ & $t_{3}$\\
$L_{d}$ & $t_{20},t_{21},t_{22}$\\
$\mathcal{A}_{0}$ & $t_{16},t_{17},t_{18},t_{19}$
\end{tabular}
\end{center}
Comparing $t_{4},\dots,t_{9}$ with $t_{12},t_{14},t_{10},t_{13},t_{15},t_{11}$ one sees that they are pairwise identical up to $L_{a}\leftrightarrow L_{b}$. Define
\begin{align}
M^{ij}\equiv\frac{1}{C}\Bigg\{\frac{c^{i}a^{j}+b^{i}c^{j}}{A-B}-\frac{1}{A-B}\left[\frac{B}{A}\,a^{i}c^{j}+\frac{A}{B}\,c^{i}b^{j}\right]
+\delta^{ij}\,\frac{B\,(c\cdot a)-A\,(c\cdot b)}{(A-B)^{2}}\Bigg\}. \label{eq:M}
\end{align}
Then $t_{4}+\dots+t_{9}=-M^{ij}L_{a}$ and $t_{10}+\dots+t_{15}=-M^{ij}L_{b}$, and the coefficients of $L_{a}$ and $L_{b}$ are \emph{the same tensor}:
\begin{align}
\mathcal{R}^{ij}=\mathcal{Q}^{ij}\,\big(L_{a}+L_{b}\big)-\frac{a^{i}b^{j}}{AB}\,L_{c}+\big(t_{20}+t_{21}+t_{22}\big)+\big(t_{16}+\dots+t_{19}\big),\qquad
\mathcal{Q}^{ij}\equiv\frac{\delta^{ij}}{2}\,\frac{A+B}{(A-B)^{2}}-M^{ij}. \label{eq:grouped}
\end{align}

\section{Step 3: the tensor $\mathcal{Q}^{ij}$ --- the poles at $A=B$ cancel}

\paragraph{$\delta^{ij}$ part.} With \eqref{eq:ids},
\begin{align}
B\,(c\cdot a)-A\,(c\cdot b)=B\,(A-a\cdot b)-A\,(a\cdot b-B)=2AB-(A+B)\,a\cdot b .
\end{align}
Hence the coefficient of $\delta^{ij}$ in $\mathcal{Q}^{ij}$ is
\begin{align}
\frac{A+B}{2(A-B)^{2}}-\frac{2AB-(A+B)\,a\cdot b}{C\,(A-B)^{2}}
&=\frac{(A+B)\,C-4AB+2(A+B)\,a\cdot b}{2\,C\,(A-B)^{2}}\notag\\
&=\frac{(A+B)(A+B-2a\cdot b)-4AB+2(A+B)\,a\cdot b}{2\,C\,(A-B)^{2}}
=\frac{(A+B)^{2}-4AB}{2\,C\,(A-B)^{2}}=\frac{1}{2C}.
\end{align}
The double pole $1/(A-B)^{2}$ has cancelled.

\paragraph{Tensor part.} Substitute $c=a-b$:
\begin{align}
c^{i}a^{j}+b^{i}c^{j}&=(a^{i}a^{j}-b^{i}a^{j})+(b^{i}a^{j}-b^{i}b^{j})=a^{i}a^{j}-b^{i}b^{j},\notag\\
\frac{B}{A}\,a^{i}c^{j}+\frac{A}{B}\,c^{i}b^{j}&=\frac{B}{A}\,(a^{i}a^{j}-a^{i}b^{j})+\frac{A}{B}\,(a^{i}b^{j}-b^{i}b^{j}).
\end{align}
The difference is
\begin{align}
a^{i}a^{j}\left(1-\frac{B}{A}\right)-b^{i}b^{j}\left(1-\frac{A}{B}\right)+a^{i}b^{j}\left(\frac{B}{A}-\frac{A}{B}\right)
=(A-B)\left[\frac{a^{i}a^{j}}{A}+\frac{b^{i}b^{j}}{B}-\frac{(A+B)\,a^{i}b^{j}}{AB}\right],
\end{align}
so the single pole $1/(A-B)$ cancels as well:
\begin{align}
M^{ij}\Big|_{\rm tensor}=\frac{1}{C}\left[\frac{a^{i}a^{j}}{A}+\frac{b^{i}b^{j}}{B}-\frac{(A+B)\,a^{i}b^{j}}{AB}\right].
\end{align}

\paragraph{Result.}
\begin{align}
\mathcal{Q}^{ij}=\frac{1}{C}\left\{\frac{\delta^{ij}}{2}-\frac{a^{i}a^{j}}{A}-\frac{b^{i}b^{j}}{B}+\frac{(A+B)\,a^{i}b^{j}}{AB}\right\}. \label{eq:Q}
\end{align}
$\mathcal{Q}^{ij}$ is regular at $(x-w)^{2}=(x-z)^{2}$.

\section{Step 4: the logarithms with sign-indefinite arguments}

\begin{center}
\renewcommand{\arraystretch}{1.4}
\begin{tabular}{c|c|c}
 & region $A>B$ & region $A<B$\\ \hline
$L_{a}=\log\dfrac{B}{B-A}$ & argument $<0$:\ \ $\log\dfrac{B}{A-B}+i\sigma\pi$ & argument $>0$:\ \ $\log\dfrac{B}{B-A}$\\
$L_{b}=\log\dfrac{A-B}{A}$ & argument $>0$:\ \ $\log\dfrac{A-B}{A}$ & argument $<0$:\ \ $\log\dfrac{B-A}{A}+i\sigma\pi$\\ \hline
$L_{a}+L_{b}$ & $\log\dfrac{B}{A}+i\sigma\pi$ & $\log\dfrac{B}{A}+i\sigma\pi$
\end{tabular}
\end{center}
In each region exactly one of the two logarithms has a negative argument, $\log|A-B|$ cancels in the sum, and
\begin{align}
L_{a}+L_{b}=L_{c}=\log\frac{B}{A}+i\,\sigma\,\pi\qquad\text{for all }A,B. \label{eq:logsum}
\end{align}
Here $\log(-|t|)=\log|t|+i\sigma\pi$ with one common $\sigma=\pm1$ for all three logarithms ($\sigma=+1$ is the principal branch); this is an assumption about the $i\epsilon$ prescription of the original integrals.

Inserting \eqref{eq:logsum} and \eqref{eq:Q} in \eqref{eq:grouped}, the three logarithms $L_{a},L_{b},L_{c}$ multiply
\begin{align}
\mathcal{Q}^{ij}-\frac{a^{i}b^{j}}{AB}
=\frac{1}{C}\left\{\frac{\delta^{ij}}{2}-\frac{a^{i}a^{j}}{A}-\frac{b^{i}b^{j}}{B}+\frac{(A+B-C)\,a^{i}b^{j}}{AB}\right\}
=\frac{1}{C}\left\{\frac{\delta^{ij}}{2}-\frac{a^{i}a^{j}}{A}-\frac{b^{i}b^{j}}{B}+\frac{2\,(a\cdot b)\,a^{i}b^{j}}{AB}\right\}\equiv K^{ij},
\end{align}
so that
\begin{align}
t_{1}+\dots+t_{15}=K^{ij}\left\{\log\frac{B}{A}+i\sigma\pi\right\}.
\end{align}

\section{Step 5: the $L_{d}$ terms}

With $c=a-b$,
\begin{align}
t_{20}+t_{21}+t_{22}=-\frac{L_{d}}{2C}\left[\frac{c^{i}b^{j}}{B}-\frac{a^{i}c^{j}}{A}+\delta^{ij}\right]
=-\frac{L_{d}}{2C}\left[\frac{a^{i}b^{j}-b^{i}b^{j}}{B}-\frac{a^{i}a^{j}-a^{i}b^{j}}{A}+\delta^{ij}\right]
=-\frac12\,L_{d}\,S^{ij},
\end{align}
\begin{align}
S^{ij}\equiv\frac{1}{C}\left\{\delta^{ij}-\frac{a^{i}a^{j}}{A}-\frac{b^{i}b^{j}}{B}+\frac{(A+B)\,a^{i}b^{j}}{AB}\right\}=\mathcal{Q}^{ij}+\frac{\delta^{ij}}{2C}=K^{ij}+\frac{\delta^{ij}}{2C}+\frac{a^{i}b^{j}}{AB}.
\end{align}

\section{Step 6: the $\mathcal{A}_{0}$ terms}

\subsection{The tensor}
\begin{align}
t_{16}+\dots+t_{19}&=\frac{\mathcal{A}_{0}}{C}\left\{c^{i}a^{j}+b^{i}c^{j}-(a\cdot b)\left[\frac{c^{i}b^{j}}{B}+\frac{a^{i}c^{j}}{A}\right]\right\}\notag\\
&=\frac{\mathcal{A}_{0}}{C}\left\{a^{i}a^{j}-b^{i}b^{j}-(a\cdot b)\left[\frac{a^{i}b^{j}-b^{i}b^{j}}{B}+\frac{a^{i}a^{j}-a^{i}b^{j}}{A}\right]\right\}\notag\\
&=\frac{\mathcal{A}_{0}}{C}\left\{\frac{A-a\cdot b}{A}\,a^{i}a^{j}+\frac{a\cdot b-B}{B}\,b^{i}b^{j}-(a\cdot b)\,\frac{A-B}{AB}\,a^{i}b^{j}\right\}
=\mathcal{A}_{0}\,T^{ij},
\end{align}
\begin{align}
T^{ij}\equiv\frac{1}{C}\left\{\frac{(c\cdot a)}{A}\,a^{i}a^{j}+\frac{(c\cdot b)}{B}\,b^{i}b^{j}-\frac{(a\cdot b)\,(A-B)}{AB}\,a^{i}b^{j}\right\}.
\end{align}

\subsection{The function $\mathcal{A}_{0}$}
$\Delta=\sqrt{AB-(a\cdot b)^{2}}=|a|\,|b|\sin\theta_{x}$ is twice the area of the triangle with corners $x,z,w$; $\theta_{x}$ is its interior angle at $x$ (between $a$ and $b$), $\theta_{z}$ the one at $z$ (between $x-z=b$ and $w-z=-c$), $\theta_{w}$ the one at $w$.
\begin{itemize}
\item $\dfrac{a\cdot b}{\Delta}=\cot\theta_{x}$, so $\arctan\dfrac{a\cdot b}{\Delta}=\dfrac{\pi}{2}-\theta_{x}$.
\item $B-a\cdot b=b\cdot(-c)=|b|\,|c|\cos\theta_{z}$ and $\Delta=|b\times c|=|b|\,|c|\sin\theta_{z}$ (because $a\times b=(b+c)\times b=c\times b$), so $\arctan\dfrac{B-a\cdot b}{\Delta}=\dfrac{\pi}{2}-\theta_{z}$.
\end{itemize}
Therefore
\begin{align}
\mathcal{A}_{0}=\frac{\pi-\theta_{x}-\theta_{z}}{\Delta}=\frac{\theta_{w}}{\Delta},\qquad
\cos\theta_{w}=\frac{(x-w)\cdot(z-w)}{|x-w|\,|z-w|},\quad0<\theta_{w}<\pi. \label{eq:A0}
\end{align}
Equivalently, since $\int_{0}^{1}\frac{dt}{Bt^{2}-2(a\cdot b)t+A}=\frac{1}{\Delta}\Big[\arctan\frac{Bt-a\cdot b}{\Delta}\Big]_{0}^{1}$,
\begin{align}
\mathcal{A}_{0}=\int_{0}^{1}\frac{dt}{(a-t\,b)^{2}}=\int_{0}^{1}\frac{dt}{\big(x_{t}-w\big)^{2}},\qquad x_{t}\equiv(1-t)\,x+t\,z, \label{eq:A0int}
\end{align}
the integral of the inverse squared distance from $w$ to the points of the segment between $x$ and $z$. $\mathcal{A}_{0}$ is real and has no logarithm.

\section{Step 7: result}

Adding Steps 4--6,
\begin{align}
\boxed{\;\mathcal{R}^{ij}=K^{ij}\left\{\log\frac{(x-z)^{2}}{(x-w)^{2}}+i\,\sigma\,\pi\right\}
-\frac12\,S^{ij}\,\log\frac{(z-w)^{2}}{(x-w)^{2}}+\mathcal{A}_{0}\,T^{ij}\;}
\end{align}
with, in the original variables,
\begin{align}
K^{ij}&=\frac{1}{(z-w)^{2}}\left\{\frac{\delta^{ij}}{2}-\frac{(x-w)^{i}(x-w)^{j}}{(x-w)^{2}}-\frac{(x-z)^{i}(x-z)^{j}}{(x-z)^{2}}+\frac{2\,(x-w)\!\cdot\!(x-z)\;(x-w)^{i}(x-z)^{j}}{(x-w)^{2}(x-z)^{2}}\right\},\notag\\
S^{ij}&=\frac{1}{(z-w)^{2}}\left\{\delta^{ij}-\frac{(x-w)^{i}(x-w)^{j}}{(x-w)^{2}}-\frac{(x-z)^{i}(x-z)^{j}}{(x-z)^{2}}+\frac{\big[(x-w)^{2}+(x-z)^{2}\big]\,(x-w)^{i}(x-z)^{j}}{(x-w)^{2}(x-z)^{2}}\right\},\notag\\
T^{ij}&=\frac{1}{(z-w)^{2}}\Bigg\{\frac{(z-w)\!\cdot\!(x-w)}{(x-w)^{2}}\,(x-w)^{i}(x-w)^{j}+\frac{(z-w)\!\cdot\!(x-z)}{(x-z)^{2}}\,(x-z)^{i}(x-z)^{j}\notag\\
&\hspace{2.2cm}-\frac{(x-w)\!\cdot\!(x-z)\,\big[(x-w)^{2}-(x-z)^{2}\big]}{(x-w)^{2}(x-z)^{2}}\,(x-w)^{i}(x-z)^{j}\Bigg\}. \label{eq:KST}
\end{align}
What happened to the 22 terms:
\begin{itemize}
\item $t_{1},\dots,t_{15}$ (three different logarithms, single and double poles at $A=B$) $\to$ one term $K^{ij}[\log(B/A)+i\sigma\pi]$, no pole;
\item $t_{20},t_{21},t_{22}$ $\to$ $-\tfrac12S^{ij}\log(C/A)$;
\item $t_{16},\dots,t_{19}$ $\to$ $\mathcal{A}_{0}T^{ij}$.
\end{itemize}
$K^{ij}$ is the bracket of Group I Row IV, Group II Row I ($1a$) and Group III Rows I and II: with $c=a-b$,
$\frac{b^{j}c^{i}}{CB}-\frac{a^{i}c^{j}}{CA}-\frac{a^{i}b^{j}}{AB}+\frac{\delta^{ij}}{2C}=K^{ij}$.

\section{Step 8: further simplification in two transverse dimensions}

Let $\theta$ be the signed angle from $a=x-w$ to $b=x-z$,
\begin{align}
\cos\theta=\frac{a\cdot b}{|a|\,|b|},\qquad\sin\theta=\frac{a\times b}{|a|\,|b|},\qquad a\times b\equiv\epsilon^{ij}a^{i}b^{j}=a^{1}b^{2}-a^{2}b^{1},\qquad\epsilon^{12}=+1,
\end{align}
and $\mathsf{R}^{ij}(\vartheta)\equiv\cos\vartheta\,\delta^{ij}+\sin\vartheta\,\epsilon^{ij}$.

\paragraph{$K^{ij}$ is a rotation.} With the reflection $\Pi^{ij}(a)\equiv\delta^{ij}-2a^{i}a^{j}/A$,
\begin{align}
\Pi^{ik}(a)\,\Pi^{kj}(b)=\delta^{ij}-\frac{2a^{i}a^{j}}{A}-\frac{2b^{i}b^{j}}{B}+\frac{4\,(a\cdot b)\,a^{i}b^{j}}{AB}=2\,C\,K^{ij},
\end{align}
and the product of two reflections is the rotation by twice the angle between them:
\begin{align}
K^{ij}=\frac{\Pi^{ik}(x-w)\,\Pi^{kj}(x-z)}{2\,(z-w)^{2}}=\frac{\mathsf{R}^{ij}(2\theta)}{2\,(z-w)^{2}},\qquad
u^{i}v^{j}K^{ij}=\frac{\cos2\theta\;u\cdot v+\sin2\theta\;u\times v}{2\,(z-w)^{2}} .
\end{align}

\paragraph{$S^{ij}$.} From $S=K+\delta/2C+a^{i}b^{j}/(AB)$ and $\tfrac12(1+\cos2\theta)=\cos^{2}\theta$, $\tfrac12\sin2\theta=\sin\theta\cos\theta$:
\begin{align}
S^{ij}=\frac{\cos\theta}{(z-w)^{2}}\,\mathsf{R}^{ij}(\theta)+\frac{(x-w)^{i}(x-z)^{j}}{(x-w)^{2}(x-z)^{2}} .
\end{align}

\paragraph{$\mathcal{A}_{0}T^{ij}$: the factor $1/\Delta$ cancels.} $T^{ij}$ vanishes when $a$ and $b$ are collinear; explicitly, with $\tilde a^{i}\equiv\epsilon^{ij}a^{j}$,
\begin{align}
T^{ij}=\frac{a\times b}{C\,AB}\Big[(a\cdot c)\,a^{i}\,\tilde b^{j}-(b\cdot c)\,\tilde a^{i}\,b^{j}\Big],
\end{align}
and since $\Delta=|a\times b|$, with the signed angle $\vartheta_{w}\equiv{\rm sgn}(a\times b)\,\theta_{w}$:
\begin{align}
\mathcal{A}_{0}\,T^{ij}=\frac{\vartheta_{w}}{(z-w)^{2}\,(x-w)^{2}\,(x-z)^{2}}\Big[(x-w)\!\cdot\!(z-w)\;(x-w)^{i}\,\widetilde{(x-z)}^{j}-(x-z)\!\cdot\!(z-w)\;\widetilde{(x-w)}^{i}\,(x-z)^{j}\Big].
\end{align}

\paragraph{Row III in two dimensions.}
\begin{align}
\boxed{\;\begin{aligned}
\mathcal{R}^{ij}&=\frac{\mathsf{R}^{ij}(2\theta)}{2\,c^{2}}\left\{\log\frac{B}{A}+i\sigma\pi\right\}
-\frac12\left[\frac{\cos\theta}{c^{2}}\,\mathsf{R}^{ij}(\theta)+\frac{a^{i}b^{j}}{AB}\right]\log\frac{c^{2}}{A}
+\frac{\vartheta_{w}}{c^{2}AB}\Big[(a\cdot c)\,a^{i}\tilde b^{j}-(b\cdot c)\,\tilde a^{i}b^{j}\Big]
\end{aligned}\;} \label{eq:2D}
\end{align}

\section{Step 9: one-parameter integral for the last two terms}

Since $(a-tb)^{2}$ equals $A$ at $t=0$ and $C$ at $t=1$,
\begin{align}
\log\frac{C}{A}=\int_{0}^{1}dt\,\frac{d}{dt}\log(a-tb)^{2}=\int_{0}^{1}dt\,\frac{2\,(tB-a\cdot b)}{(a-tb)^{2}},
\end{align}
and with \eqref{eq:A0int} the last two terms of \eqref{eq:main} combine to
\begin{align}
-\frac12\,S^{ij}\log\frac{C}{A}+\mathcal{A}_{0}T^{ij}&=\int_{0}^{1}dt\;\frac{T^{ij}-(tB-a\cdot b)\,S^{ij}}{(a-tb)^{2}}
=\int_{0}^{1}dt\;\frac{\mathcal{N}^{ij}(t)}{(z-w)^{2}\,\big(x_{t}-w\big)^{2}}, \label{eq:irep}\\
\mathcal{N}^{ij}(t)&=-(tB-a\cdot b)\,\delta^{ij}+\frac{A-2a\cdot b+tB}{A}\,a^{i}a^{j}-(1-t)\,b^{i}b^{j}-\frac{t(A+B)-2a\cdot b}{A}\,a^{i}b^{j},
\end{align}
where $x_{t}=(1-t)x+tz$. (The coefficients follow from $c\cdot a+tB-a\cdot b=A-2a\cdot b+tB$, $c\cdot b+tB-a\cdot b=-(1-t)B$, and $(a\cdot b)(A-B)+(tB-a\cdot b)(A+B)=B\,[t(A+B)-2a\cdot b]$.)

\section{Properties of the result}

\begin{enumerate}
\item \textbf{No singularity at $(x-w)^{2}=(x-z)^{2}$.} All poles of the original terms have cancelled and $\log(B/A)$ is regular there.
\item \textbf{Collinear configurations.} $\mathcal{A}_{0}=\theta_{w}/\Delta$ diverges when $w$ lies on the segment between $x$ and $z$, but $T^{ij}\propto a\times b$ vanishes there and $\mathcal{A}_{0}T^{ij}$ stays bounded (it changes sign across the segment through $\vartheta_{w}$).
\item \textbf{The imaginary constant.} The only remnant of the negative arguments is $i\sigma\pi K^{ij}$. With the prefactor $i$ of \eqref{eq:def} it is the real kernel $-\sigma\pi K^{ij}$:
\begin{align}
\frac{d^{3}N_{\rm I.III}}{d^{3}k}\Bigg|_{\pi}=-\,\sigma\,\frac{1}{(2\pi)^{3}}\,\frac{g^{4}}{8\pi^{4}}\,\frac{1}{k^{+}}\int e^{-ik\cdot(w'-w)}\;u^{i}\,v^{j}\,K^{ij}\;\mathcal{W}+{\rm c.c.}
\end{align}
\item \textbf{Reality.} Row III is $X+{\rm c.c.}$ and therefore real. For commuting $\rho$, $\mathcal{W}$ is real and
\begin{align}
\frac{d^{3}N_{\rm I.III}}{d^{3}k}&=\frac{1}{(2\pi)^{3}}\,\frac{g^{4}}{4\pi^{5}}\,\frac{1}{k^{+}}\int u^{i}v^{j}\,\mathcal{W}
\Big\{\sin\!\big[k\cdot(w'-w)\big]\Big(K^{ij}\log\frac{B}{A}-\tfrac12S^{ij}\log\frac{C}{A}+\mathcal{A}_{0}T^{ij}\Big)\notag\\
&\hspace{4.2cm}-\sigma\pi\,\cos\!\big[k\cdot(w'-w)\big]\,K^{ij}\Big\}.
\end{align}
\item \textbf{Behaviour at $w\to z$ ($c\to0$).} $K^{ij}\to\delta^{ij}/2C$, $S^{ij}\to\delta^{ij}/C$, while $K^{ij}\log(B/A)$ and $\mathcal{A}_{0}T^{ij}$ are $\mathcal{O}(1/|c|)$:
\begin{align}
\mathcal{R}^{ij}\;\longrightarrow\;-\frac{\delta^{ij}}{2\,(z-w)^{2}}\left\{\log\frac{(z-w)^{2}}{(x-w)^{2}}-i\sigma\pi\right\}.
\end{align}
$\mathcal{W}$ does not depend on $w$ and $\int d^{2}w\,\log(c^{2})/c^{2}$ does not converge at $w\to z$. The source is $t_{22}$ (the last term of $N_{3b}$) and the $i\sigma\pi$ term. So Row III as given is not finite by itself at $w\to z$.
\item \textbf{Relation to the other rows.} The piece $K^{ij}\log(B/A)$ has the same kernel as Group I Rows II and IV, Group II Row I (first part) and Group III Rows I, II; together these add up to commutators of $\rho$ only (see the general note). What is specific to Row III is $i\sigma\pi K^{ij}-\tfrac12S^{ij}\log(C/A)+\mathcal{A}_{0}T^{ij}$.
\end{enumerate}

\section{Numerical checks}
\begin{center}
\renewcommand{\arraystretch}{1.2}
\begin{tabular}{p{9cm}|p{7.5cm}}
statement & check\\ \hline
$\sum_{n=1}^{22}t_{n}=K L_{c}-\tfrac12SL_{d}+\mathcal{A}_{0}T$ & 1500 random $(a,b)$, complex principal-branch logs, relative deviation $<10^{-10}$\\
coefficients of $L_{a}$ and of $L_{b}$ both equal $\mathcal{Q}^{ij}$ & same run, $10^{-11}$\\
$\mathcal{A}_{0}=\theta_{w}/\Delta=\int_{0}^{1}dt/(a-tb)^{2}$ & same run; quadrature, $10^{-16}$\\
two-dimensional forms of $K$, $S$, $\mathcal{A}_{0}T$ & same run, $10^{-9}$ or better\\
integral representation \eqref{eq:irep} & quadrature, $10^{-15}$
\end{tabular}
\end{center}

\end{document}
